% Phụ lục B

\chapter{Liệt kê source code} % Tên của phụ lục

\label{AppendixB} % Để trích dẫn chương này ở chỗ nào đó trong bài, hãy sử dụng lệnh \ref{AppendixB} 

\begin{lstlisting}[language=Python, caption={Object Detection Function}, label={lst:detect_function}]
def detect(self, frame: np.ndarray) -> List[Dict[str, Any]]:
    if self.model is None:
        return self._mock_detect(frame)

    try:
        results = self.model(
            frame,
            conf=self.conf,
            iou=self.iou,
            imgsz=self.imgsz,
            verbose=False,
        )[0]

        detections = []
        for box in results.boxes:
            cls_id = int(box.cls[0])
            cls_name = results.names.get(cls_id, "default")
            conf = float(box.conf[0])
            x1, y1, x2, y2 = [int(v) for v in box.xyxy[0]]
            cx = (x1 + x2) // 2
            cy = (y1 + y2) // 2

            detections.append({
                "class_id": cls_id,
                "class_name": cls_name,
                "label_vi": VI_LABELS.get(cls_name, VI_LABELS["default"]),
                "confidence": round(conf, 3),
                "bbox": [x1, y1, x2, y2],
                "center": [cx, cy],
                "area": (x2 - x1) * (y2 - y1),
            })
        return detections

    except Exception as e:
        logger.error(f"Detection error: {e}")
        return []
\end{lstlisting}

\begin{lstlisting}[language=Python, caption={Object Tracking Function}, label={lst:tracking_update}]
def update(self, detections: List[Dict[str, Any]]) -> List[Dict[str, Any]]:
    if not detections:
        for obj_id in list(self.disappeared.keys()):
            self.disappeared[obj_id] += 1
            if self.disappeared[obj_id] > self.max_disappeared:
                self._deregister(obj_id)
        return []

    input_centroids = [d["center"] for d in detections]

    if not self.objects:
        for i, det in enumerate(detections):
            self._register(det)
    else:
        object_ids = list(self.objects.keys())
        object_centroids = [
            self.objects[oid]["center"] for oid in object_ids
        ]

        # Simple nearest-neighbour matching
        matched_det = set()
        matched_obj = set()

        for i, obj_id in enumerate(object_ids):
            min_dist = float("inf")
            min_j = -1

            for j, inp_c in enumerate(input_centroids):
                if j in matched_det:
                    continue

                dist = self._euclidean(
                    object_centroids[i], inp_c
                )

                if dist < min_dist:
                    min_dist = dist
                    min_j = j

            if min_j >= 0 and min_dist < self.max_distance:
                self.objects[obj_id] = {
                    **detections[min_j],
                    "track_id": obj_id,
                }

                self.history[obj_id].append(
                    detections[min_j]["center"]
                )

                if len(self.history[obj_id]) > 15:
                    self.history[obj_id] = \
                        self.history[obj_id][-15:]

                self.disappeared[obj_id] = 0
                matched_det.add(min_j)
                matched_obj.add(obj_id)

            else:
                self.disappeared[obj_id] += 1

                if (
                    self.disappeared[obj_id]
                    > self.max_disappeared
                ):
                    self._deregister(obj_id)

        for j, det in enumerate(detections):
            if j not in matched_det:
                self._register(det)

    # Enrich with motion info
    result = []

    for obj_id, obj in self.objects.items():
        enriched = {**obj}
        enriched["track_id"] = obj_id

        (
            enriched["direction"],
            enriched["speed_label"],
            enriched["speed_score"],
        ) = self._analyze_motion(obj_id)

        result.append(enriched)

    return result
\end{lstlisting}

\begin{lstlisting}[language=Python, caption={Depth Estimation Function Using MiDaS}, label={lst:depth_estimation}]
def estimate(self, frame: np.ndarray) -> np.ndarray:
    """Returns normalized depth map [0..1], 1=near, 0=far"""
    
    if self.model is None:
        return self._gradient_fallback(frame)

    try:
        import torch

        rgb = cv2.cvtColor(frame, cv2.COLOR_BGR2RGB)
        inp = self.transform(rgb).to(self.device)

        with torch.no_grad():
            pred = self.model(inp)

            pred = torch.nn.functional.interpolate(
                pred.unsqueeze(1),
                size=frame.shape[:2],
                mode="bicubic",
                align_corners=False,
            ).squeeze()

        depth = pred.cpu().numpy()

        # Invert: MiDaS output is disparity (higher = closer)
        d_min, d_max = depth.min(), depth.max()

        if d_max > d_min:
            depth = (depth - d_min) / (d_max - d_min)

        return depth

    except Exception as e:
        logger.error(f"Depth estimation error: {e}")
        return self._gradient_fallback(frame)
\end{lstlisting}

\begin{lstlisting}[language=Python, caption={Text-to-Speech Audio Generation Functions}, label={lst:tts_functions}]
def _gtts(self, text: str, fpath: Path) -> bool:
    try:
        from gtts import gTTS

        tts = gTTS(
            text=text,
            lang=self.lang,
            slow=False
        )

        tts.save(str(fpath))

        ok = (
            fpath.exists()
            and fpath.stat().st_size > 1000
        )

        if not ok:
            fpath.unlink(missing_ok=True)

        return ok

    except Exception as e:
        logger.debug(f"gTTS failed: {e}")
        fpath.unlink(missing_ok=True)
        return False


def _espeak(self, text: str, fpath: Path) -> bool:
    try:
        import pyttsx3

        engine = pyttsx3.init()

        voices = engine.getProperty("voices")
        vi_voice = next(
            (
                v for v in voices
                if "vi" in v.id.lower()
            ),
            None,
        )

        if vi_voice:
            engine.setProperty(
                "voice",
                vi_voice.id
            )

        engine.setProperty("rate", 145)
        engine.setProperty("volume", 1.0)

        wav = str(fpath).replace(
            ".mp3",
            ".wav"
        )

        engine.save_to_file(text, wav)
        engine.runAndWait()
        engine.stop()

        if (
            os.path.exists(wav)
            and os.path.getsize(wav) > 500
        ):

            # Try ffmpeg wav -> mp3
            ret = os.system(
                f'ffmpeg -y -i "{wav}" '
                f'-ar 22050 -ac 1 -q:a 4 '
                f'"{fpath}" '
                f'-loglevel quiet '
                f'>nul 2>&1'
            )

            os.remove(wav)

            if not (
                fpath.exists()
                and fpath.stat().st_size > 500
            ):

                # ffmpeg failed
                engine2 = pyttsx3.init()

                if vi_voice:
                    engine2.setProperty(
                        "voice",
                        vi_voice.id
                    )

                engine2.setProperty(
                    "rate",
                    145
                )

                engine2.save_to_file(
                    text,
                    str(fpath)
                )

                engine2.runAndWait()
                engine2.stop()

            return (
                fpath.exists()
                and fpath.stat().st_size > 500
            )

        return False

    except Exception as e:
        logger.debug(f"eSpeak failed: {e}")
        return False
\end{lstlisting}